%! Sinusoidal Function Transformations — Unit 3, Lesson 3.6 — Homework
\documentclass[10pt]{article}
\usepackage{precalculus-article}
\usepackage{precalculus-boxes}
\newcommand{\TallMath}[1]{$\displaystyle #1\rule[-1.4em]{0pt}{3.2em}$}

\begin{document}

\pageheader{Unit 3, Lesson 3.6}{Homework}
\namedateperiod

\vspace{0.10in}

\begin{notesbox}{Problems 1--6}

\textbf{1.}\quad For each function, identify the amplitude and midline.

\vspace{4pt}
\renewcommand{\arraystretch}{1.5}
\begin{tabular}{|l|c|c|}
\hline
\textbf{Function} & \textbf{Amplitude} & \textbf{Midline} \\
\hline
$f(\theta) = 7\sin(\theta) + 3$ & & \\
\hline
$g(\theta) = -5\cos(\theta) - 4$ & & \\
\hline
$h(\theta) = \tfrac{3}{2}\sin(2\theta) + 1$ & & \\
\hline
$k(\theta) = -\cos(\pi\theta)$ & & \\
\hline
\end{tabular}

\vspace{0.12in}
\textbf{2.}\quad For each function, identify the period by first finding $b$.

\vspace{4pt}
(a)\quad $f(\theta) = \sin(4\theta)$ \qquad $b = $ \underline{\hspace{1.2cm}}
\qquad Period $= \dfrac{2\pi}{|b|} = $ \underline{\hspace{2cm}}

\vspace{6pt}
(b)\quad $g(\theta) = \cos\!\left(\dfrac{\pi}{3}\theta\right)$ \qquad $b = $ \underline{\hspace{1.2cm}}
\qquad Period $= $ \underline{\hspace{2cm}}

\vspace{6pt}
(c)\quad $h(\theta) = 3\sin\!\left(\dfrac{1}{4}\theta\right) - 2$ \qquad Period $= $ \underline{\hspace{2cm}}

\vspace{0.12in}
\textbf{3.}\quad For each function, factor the argument and identify the phase shift.
State the direction and amount.

\vspace{4pt}
(a)\quad $f(\theta) = \sin(2\theta - \pi)$

\vspace{2pt}
Factor: $2\theta - \pi = 2\!\left(\theta - \underline{\hspace{1.5cm}}\right)$
\quad Phase shift: \underline{\hspace{3.5cm}}

\vspace{6pt}
(b)\quad $g(\theta) = \cos\!\left(3\theta + \dfrac{\pi}{2}\right)$

\vspace{2pt}
Factor: $3\theta + \dfrac{\pi}{2} = 3\!\left(\theta + \underline{\hspace{1.5cm}}\right)$
\quad Phase shift: \underline{\hspace{3.5cm}}

\vspace{6pt}
(c)\quad $h(\theta) = \sin\!\left(\dfrac{1}{2}\theta - \dfrac{\pi}{4}\right)$

\vspace{2pt}
Factor: \underline{\hspace{4cm}}
\quad Phase shift: \underline{\hspace{3.5cm}}

\vspace{0.12in}
\textbf{4.}\quad Identify all four features (amplitude, period, phase shift, midline)
for each function.

\vspace{4pt}
(a)\quad $f(\theta) = 4\sin\!\left(3\!\left(\theta - \dfrac{\pi}{6}\right)\right) + 5$

\vspace{4pt}
Amplitude: \underline{\hspace{1.8cm}} \quad
Period: \underline{\hspace{1.8cm}} \quad
Phase shift: \underline{\hspace{2.2cm}} \quad
Midline: \underline{\hspace{1.8cm}}

\vspace{8pt}
(b)\quad $g(\theta) = -2\cos(4\theta + \pi) - 3$

\vspace{4pt}
First factor: $4\theta + \pi = 4\!\left(\theta + \underline{\hspace{1.2cm}}\right)$

\vspace{4pt}
Amplitude: \underline{\hspace{1.8cm}} \quad
Period: \underline{\hspace{1.8cm}} \quad
Phase shift: \underline{\hspace{2.2cm}} \quad
Midline: \underline{\hspace{1.8cm}}

\vspace{0.12in}
\textbf{5. (AP-Style MC)}\quad Which of the following is the period of
$h(\theta) = 5\sin(3\theta - \pi) + 2$?

\begin{enumerate}[label=(\Alph*), leftmargin=2em, itemsep=3pt]
  \item $\pi/3$
  \item $2\pi/3$
  \item $3\pi$
  \item $\pi$
  \item $2\pi$
\end{enumerate}

Answer: \underline{\hspace{1.5cm}}

\vspace{0.12in}
\textbf{6. (AP-Style MC)}\quad The function $f(\theta) = a\sin(b\theta) + d$
has amplitude 6 and midline $y = -1$. Which of the following could be $f$?

\begin{enumerate}[label=(\Alph*), leftmargin=2em, itemsep=3pt]
  \item $f(\theta) = 6\sin(\theta) + 1$
  \item $f(\theta) = -6\sin(\theta) - 1$
  \item $f(\theta) = 6\sin(\theta) - 7$
  \item $f(\theta) = 3\sin(2\theta) - 1$
  \item $f(\theta) = 6\sin(\theta) + 7$
\end{enumerate}

Answer: \underline{\hspace{1.5cm}}

\end{notesbox}

\vspace{0.08in}

\begin{extensionbox}
\textbf{Extension (Optional):} The function $p(\theta) = A\cos(B\theta) + D$
has maximum value 7 and minimum value $-1$.

\vspace{4pt}
(a)\quad Find the amplitude $A$ and midline $D$ of $p$.

\vspace{4pt}
Amplitude $= \dfrac{\text{max} - \text{min}}{2} = $ \underline{\hspace{3cm}}
\quad Midline $D = \dfrac{\text{max} + \text{min}}{2} = $ \underline{\hspace{3cm}}

\vspace{4pt}
(b)\quad If the period of $p$ is $4\pi$, find $B$.

\vspace{4pt}
(c)\quad Rewrite $p(\theta)$ as a sine function with an appropriate phase shift.
Show all work.
\writelines{3}
\end{extensionbox}

\vspace{0.08in}

\begin{tcolorbox}[colback=navylight, colframe=navy, arc=2mm,
    left=3mm, right=3mm, top=2mm, bottom=2mm]
\small\textbf{Preview --- Lesson 3.7: Sinusoidal Function Graphs}\\
You can now read amplitude, period, phase shift, and midline directly from
a sinusoidal equation. In Lesson 3.7 we'll go in both directions: from
equation to graph (using key points), and from a graph back to the equation.
Come ready to use today's parameters to plot and interpret complete sinusoidal graphs.
\end{tcolorbox}

\end{document}
